\scanpage{193}
Les $\unsure{\xi}_l$ de (i) dépendent du choix d'un pt géom de $X$,
mais à isom non unique près (dépendant du
\uuline{$\pi_1(X)$}) \inserted{[$\pi_1(X)$ opère sur $\unsure{\xi}_l$]}
\deleted{sont} \unsure{indépendants} de ce choix.

Les $\unsure{\xi}_p^\alpha$ dépendent du choix d'un \deleted{corps} pt\inserted{de $X$}
à valeurs dans un corps parfait sans plus,
si on veut les $k(x)^{\mathrm{pf}}$ (clôture parfaite de $k(x)$),
et pour $k$ donné, il ne \uline{semble pas} qu'il y ait
\deleted{ait} des isom entre les $\unsure{\xi}_{p,a}^{(\alpha)}$ correspondants
aux divers pts à valeurs dans $k$ (c'est la
``partie continue'' du $\unsure{H^1|H^2}$ préc. qui v. faire des
ennuis), s'ils sont localisés en le $\overline{\mathrm m}$ $x\in X$, \deleted{il}\inserted{et}
\deleted{sont} si $k$ est alg. clos, ils sont conjugués par
autom. de $k$\inserted{(étendus à $\mathbb W_p(k)$)}, les groupes alg. associés sont
donc conjugués par autom. de $k$ $-\, -\, -\, -$ il est $\overline{\mathrm m}$
gentil qu'ils \unsure{veulent|veuillent} être l'un \unsure{tordu} de
l'autre $-\, -\, -\, -$ Le gpe $\pi_1(X)$ n'opère pas sur $\unsure{\xi}_p^\alpha$
\figure{\begin{tikzpicture}[baseline=-.4ex,scale=.22]\draw (0,0)--(1,2);\draw (1,2) circle (.18);\end{tikzpicture}}, par
contre le groupe des autom de $k$ sur $k(x)$
\deleted{[\unsure{au $\overline{\mathrm m}$ de $k$} \unsure{$H$ précédent}]}
opère de façon $1/2$ \uline{linéaire} sur $\unsure{\xi}_p^{(\alpha)}$ $-\, -\, -\, -$
\margin{[\unsure{En une descente}
\figure{\begin{tikzpicture}[baseline=-.4ex,scale=.22]\draw (0,0)--(1,2);\draw (0,0) circle (.18);\end{tikzpicture}}}

\margin{\figure{\begin{tikzpicture}[baseline=-.4ex,scale=.22]\draw (0,0)--(1,2);\draw (1,2) circle (.18);\end{tikzpicture}}
En revanche, on a comme structure supplémentaire
les opérations $\unsure{F\,V|F,V}$ avec $FV=p$}

\margin{\figure{\begin{tikzpicture}[baseline=-.4ex,scale=.22]\draw (0,0)--(1,2);\draw (0,0) circle (.18);\end{tikzpicture}}
\unsure{élégante.} D'autre part, un autom $\sigma$
de $k$ qui n'est pas $k(x)$-linéaire
transforme en général $\unsure{\xi}_{xp}^{\unsure{\alpha}}$ en un
$\unsure{\xi}_{xp}^{\sigma}$ qui n'est plus isom
à $\unsure{\xi}_{x,p}$ $-\, -\, -\, -$}

Notations

\diagram{\[
\text{(i)}\left\{\begin{array}{ll}
\unsure{\xi}_{\bar x,l} & \text{foncteurs fibres $l$-adiques en}\\
& \leftarrow\ \text{pt géom }\bar x\\[6pt]
\unsure{\xi}_{x,l} & \text{par desc. de la \unsure{$2$-gpe|$l$-gpe} pour $x\in X$}\\
& \text{ou un pt de $X$ à val.}\\
& \text{dans un corps quelc.}\\
& \text{($\unsure{\xi}_{x,l}$ déf. à isom \uline{non uniq.} près)}
\end{array}\right.
\]}

\margin{\uline{On peut mettre}\\
\uline{des \unsure{exposants} $\alpha$}\\
\uline{à \unsure{volonté}}}

\diagram{\[
\text{(ii)}\left\{\begin{array}{ll}
\unsure{\xi}_{x,p} & \text{$x$ \unsure{est pt} de $X$ à val. dans un corps}\\[6pt]
\unsure{\xi}_{x,p}^{0} &
\end{array}\right.
\]}

\margin{192}
